\section{Curvature of a varying stitch length}
\label{sec:h-gravw}

A localized excess or deficit of stitches changes \(\ell_W(x)\) and
therefore the transverse coefficient in
\begin{equation}
\label{eq:dsx}
ds^2=dx^2+dy^2+\ell_W(x)^2\,dz^2.
\end{equation}
The curvature of a diagonal metric is a function of \(\ell_W(x)\) and its
derivatives.
The background itself is anisotropic by \(\ell_W\), and on the woven
ladder \(\ell_W\ge 1.77\) (Proposition~\ref{prop:over}).
A linearization about \(\delta_{ab}\) is a different metric.

Perturbations on \eqref{eq:dsx} fall into two classes.
At fixed \(\ell_W\), the in-plane static response is the Helmholtz problem
\eqref{eq:helm}, or the excluded Laplace problem at the band edge
(Section~\ref{sec:h-grav}).
The range is \(O(1)\), or \(1/r\) only on that excluded tuning.
Gradients of \(\ell_W\) are changes of \(G\).
They are absent between events (Section~\ref{sec:h-gravq}).

If \(\ell_W\) is uniform, the curvature associated with its derivatives
vanishes and a packet propagates on the constant anisotropic background.
A quadrupole aligned with the sheet normal, scaling as the second
derivatives of \(\ell_W\), is present when the stitch field varies in
space.
An isotropic \(1/r\) potential is not a solution of \eqref{eq:dsx}.
Section~\ref{sec:h-metric} is the same transverse coefficient read from
the Hessian rather than from the curvature.
